\subsection*{Bijective Functions}
\label{subsec:bijective-functions}

\begin{definitionbox}{Definition (Bijective)}
\begin{definition}[Bijective Function]
\label{def:alg-bijective}
Let $f:A\to B$.
The function $f$ is \textbf{bijective} if it is both injective and surjective.
\end{definition}
\end{definitionbox}

\begin{remark*}[Standard quantified statement]
\[
\operatorname{Bijective}(f,A,B)
\Longleftrightarrow
\operatorname{Injective}(f,A,B)\land \operatorname{Surjective}(f,A,B).
\]
\end{remark*}

\begin{remark*}[Definition predicate reading]
$f$ is bijective exactly when it is both injective and surjective.
\end{remark*}

\begin{remark*}[Negated quantified statement]
\[
\neg\operatorname{Bijective}(f,A,B)
\Longleftrightarrow
\neg\operatorname{Injective}(f,A,B)
\lor
\neg\operatorname{Surjective}(f,A,B).
\]
\end{remark*}

\begin{remark*}[Interpretation]
Bijectivity is the set-theoretic part of invertibility.  A bijective
structure-preserving map is the usual route to an isomorphism.
\end{remark*}

\begin{dependencies}
\begin{itemize}
  \item \hyperref[def:alg-injective]{Injective Function}
  \item \hyperref[def:alg-surjective]{Surjective Function}
\end{itemize}
\end{dependencies}
